\documentclass[11pt]{article}

\usepackage[T1]{fontenc}
\usepackage[utf8]{inputenc}
\usepackage{lmodern}
\usepackage[a4paper,margin=1in]{geometry}
\usepackage{amsmath,amssymb,amsthm,mathtools}
\usepackage{booktabs}
\usepackage{array}
\usepackage{xcolor}
\usepackage{microtype}
\usepackage{enumitem}
\usepackage[hidelinks]{hyperref}

\setlength{\parindent}{0pt}
\setlength{\parskip}{0.55em}
\setlength{\emergencystretch}{3em}

\newtheorem{lemma}{Lemma}[section]
\newtheorem{theorem}[lemma]{Theorem}
\newtheorem{proposition}[lemma]{Proposition}
\newtheorem{invariant}[lemma]{Invariant}
\theoremstyle{definition}
\newtheorem{definition}[lemma]{Definition}
\newtheorem{example}[lemma]{Example}
\theoremstyle{remark}
\newtheorem{remark}[lemma]{Remark}

\newcommand{\Tid}{\mathsf{Tid}}
\newcommand{\Loc}{\mathsf{Loc}}
\newcommand{\Val}{\mathsf{Val}}
\newcommand{\Node}{\mathsf{Node}}
\newcommand{\Pos}{\mathsf{Pos}}
\newcommand{\dom}{\mathop{\mathrm{dom}}}
\newcommand{\gap}{\mathop{\mathrm{gap}}}
\newcommand{\rank}{\mathop{\mathrm{rank}}}
\newcommand{\Ready}{\mathsf{Ready}}
\newcommand{\Atomic}{\mathsf{Atomic}}
\newcommand{\Inv}{\mathsf{inv}}
\newcommand{\Lini}{\mathsf{lini}}
\newcommand{\Linr}{\mathsf{linr}}
\newcommand{\Push}{\mathsf{push}}
\newcommand{\TryPop}{\mathsf{trypop}}
\newcommand{\Succ}{\mathsf{Succ}}
\newcommand{\Fail}{\mathsf{Fail}}
\newcommand{\Empty}{\mathsf{Empty}}
\newcommand{\true}{\mathsf{true}}
\newcommand{\false}{\mathsf{false}}
\newcommand{\poss}{\Delta}
\newcommand{\conc}{\sigma}
\newcommand{\abst}{\rho}
\newcommand{\tokens}{\pi}
\newcommand{\entails}{\mathrel{\vDash}}
\newcommand{\code}[1]{\texttt{\detokenize{#1}}}
\newcommand{\now}{\mathsf{now}}
\newcommand{\Real}{\mathsf{Real}}
\newcommand{\Rep}{\mathsf{Rep}}
\newcommand{\Tok}{\mathsf{Tok}}
\newcommand{\Placed}{\mathsf{placed}}
\newcommand{\Popped}{\mathsf{popped}}
\newcommand{\RetBefore}{\mathsf{returned}}
\newcommand{\Blk}{\mathsf{Blk}}
\newcommand{\FS}{\mathsf{FS}}
\newcommand{\Safe}{\mathsf{Safe}}
\newcommand{\Newer}{\mathsf{newer}}
\newcommand{\cut}{\mathsf{cut}}
\newcommand{\rem}{\mathsf{rem}}
\newcommand{\snap}{\mathsf{snap}}
\newcommand{\phase}[1]{\mathsf{#1}}
\newcommand{\PhaseIs}{\mathsf{PhaseIs}}
\newcommand{\Done}{\mathsf{Done}}
\newcommand{\Inside}{\mathsf{Inside}}
\newcommand{\lt}{\prec}
\newcommand{\lte}{\preceq}
\newcommand{\ret}{\mathsf{ret}}
\newcommand{\pop}{\mathsf{pop}}
\newcommand{\inv}{\mathsf{inv}}
\newcommand{\steps}{\mathsf{steps}}

\title{The TryStack Layer: A Detailed Account of the Mechanized Proof\\
  \large from the snapshot-based TryStackAux to the atomic TryStack}
\author{Proof exposition for the mechanized development}
\date{\today}

\begin{document}
\maketitle

\begin{abstract}
This note explains the correctness proof of the TryStack layer.  The
implementation is an identity adapter: it forwards \(\Push\) and \(\TryPop\)
to the TryStackAux underlay.  All the difficulty is in the specifications.
TryStackAux performs a \(\TryPop\) by taking a snapshot of the vertex set at
invocation and removing, at response time, a vertex that is top \emph{in that
snapshot}; TryStack pops a top of the \emph{current} live graph atomically.
Linearizing the former as the latter requires the abstract pop to be placed
retroactively, sometimes before the abstract linearization of a push that was
already visible concretely.  The proof therefore keeps, for every reachable
concrete state, the set of \emph{all} abstract configurations that can be
realized by a consistent placement of abstract events into the concrete
timeline.  We present the ghost-timed concrete state, the notion of an
abstract history and its consistency conditions, the representation of an
abstract configuration at a cut, the invariant, the rely--guarantee relations,
the replay argument that simulates each concrete event, the trace theorem
that guarantees the invariant is never vacuous, and the assembly of the
method triples.  Coq-specific engineering is mentioned only where it reflects
a genuine mathematical point.
\end{abstract}

\tableofcontents

\section{Statement of the refinement}

Fix a thread domain \(\Upsilon\subseteq\Tid\) and a value type \(\Val\).  A
node identifier is a pair
\[
  n=(\beta,\ell)\in\Node\triangleq\Tid\times\Loc,
\]
where \(\beta\) is the thread that pushed the node.  The layer imports the
TryStackAux interface \(\mathcal E\) and exports the TryStack interface
\(\mathcal F\); both offer the same two operations
\[
  \Push(v)\ :\ \mathbf 1,\qquad
  \TryPop\ :\ \{\Succ(v,\beta,\ell),\ \Succ(\Empty),\ \Fail\}.
\]
The implementation is the identity adapter
\[
  \Push_\alpha(v):\ \ r\leftarrow\mathcal E.\Push(v);\ \mathsf{return}(),
  \qquad
  \TryPop_\alpha:\ \ r\leftarrow\mathcal E.\TryPop;\ \mathsf{return}(r).
\]
The results are
\begin{itemize}[nosep]
\item \(\mathsf{MTryStack}\): a layer implementation simulation
  \(\mathcal E\to\mathcal F\) obtained from the soundness theorem of the
  set-valued rely--guarantee logic;
\item \(\mathsf{MTryStackLinearizable}\): the corresponding linearizability
  statement;
\item \(\mathsf{MListPoolTryStack}\): its vertical composition with the
  ListPool-to-TryStackAux proof, a linearizability statement from the
  ListPool underlay to the atomic TryStack.
\end{itemize}
The development has no admitted facts.  Beyond classical logic and the
extensionality principles already used by the framework, no axiom is needed.

\section{The two specifications}

Both specifications describe the stack as a DAG whose vertices are pushed
nodes.  An edge \(x\to y\) records that the push of \(x\) started after the
push of \(y\) had completed; two overlapping pushes are unordered.  A
\emph{top} vertex is a live vertex with no live vertex above it.

\subsection{The atomic TryStack}

A TryStack state is a tuple \(\abst=(V,E,P,g)\): the vertex map
\(V:\Node\rightharpoonup\Val\), the edge relation \(E\), the pending-push map
\(P:\Tid\rightharpoonup\Loc\), and the garbage set \(g\).  A vertex is live if
it is defined and not garbage.  The transitions are:
\begin{itemize}[nosep]
\item \(\Push(v)\) invocation by \(\alpha\) at a fresh location \(\ell\):
  add \(V(\alpha,\ell)=v\), add edges from \((\alpha,\ell)\) to every
  defined, non-pending vertex, set \(P(\alpha)=\ell\).
\item \(\Push\) response: remove \(P(\alpha)\).
\item \(\TryPop\): an invocation moves to an atomic-pending control state;
  the response either marks a top live vertex \(n\) as garbage and returns
  \(\Succ(V(n),n)\), returns \(\Succ(\Empty)\) when every vertex is garbage,
  or returns \(\Fail\) without changing the state.
\end{itemize}
Push is interval-sequential (its two events may be separated by other
threads' events); try-pop is atomic (its two events are adjacent, enforced by
the control state \(\Atomic(\abst,\alpha,\TryPop)\)).

\subsection{The snapshot-based TryStackAux}

A TryStackAux state is \(\conc=(V,E,S,P,g)\) with an additional snapshot map
\(S:\Tid\rightharpoonup\mathcal P(\Node)\).  Push behaves as above.  A
\(\TryPop\) by \(\alpha\) has two branches:
\begin{itemize}[nosep]
\item \emph{snapshot branch}: the invocation records \(S(\alpha)=\dom V\);
  the response either removes a vertex \(n\) that is top in
  \(S(\alpha)\setminus g\) (with respect to \(E\)) and returns
  \(\Succ(V(n),n)\), or returns \(\Fail\); both clear \(S(\alpha)\);
\item \emph{empty branch}: when every vertex is garbage the invocation may
  enter an atomic-pending state whose only response is \(\Succ(\Empty)\).
\end{itemize}
The crucial difference is that the removed vertex \(n\) is top only among
the vertices that existed at the snapshot.  A vertex \(x\) pushed after the
snapshot may already sit above \(n\) in the current graph.

\subsection{Ghost timing}

To reason about which retroactive placements are possible, the TryStackAux
state carries ghost fields that never restrict a transition:
\begin{itemize}[nosep]
\item a logical clock \(\now\), advanced by every recorded event;
\item \(i_n\), the clock value at which the push of \(n\) was invoked, and
  \(r_n\), the clock value of its response (undefined while pending);
\item \(\snap(\alpha)\), the clock value of \(\alpha\)'s pending snapshot;
\item \(\rem(n)=(\alpha,st,rt)\) for a removed vertex: the remover, the
  clock value of its snapshot and the clock value of the removal.
\end{itemize}
Every recorded clock value is distinct.  We write \(x\to_\conc y\) for the
concrete edge, which the well-formedness invariant identifies with
\(r_y<i_x\).

\begin{invariant}[Ghost well-formedness, \code{tsa_ghost_wf}]
\label{inv:wf}
The following hold in every reachable state and are preserved by every step.
\begin{enumerate}[nosep]
\item \(n\) is a vertex iff \(i_n\) is defined; \(i_n<\now\); if \(r_n\) is
  defined then \(i_n<r_n<\now\); \(n\) is pending iff \(i_n\) is defined
  and \(r_n\) is not.
\item \(x\to_\conc y\iff r_y<i_x\).
\item \(S(\alpha)=N\) iff \(\snap(\alpha)=t\) is defined and
  \(N=\{n\mid i_n<t\}\).
\item \(n\in g\) iff \(\rem(n)\) is defined.
\item (\emph{Blocking}) If \(\rem(n)=(\alpha,st,rt)\) then
  \(i_n<st<rt<\now\), and every \(y\) with \(i_y<st\) and \(y\to_\conc n\)
  satisfies \(\rem(y)=(\_,\_,rt')\) with \(rt'<rt\).
\item Clock values are unique across all recorded events.
\item A thread pushes sequentially: if \(i_{(\alpha,\ell_1)}<i_{(\alpha,\ell_2)}\)
  then \(r_{(\alpha,\ell_1)}\) is defined and below \(i_{(\alpha,\ell_2)}\).
\end{enumerate}
\end{invariant}

Clause~5 is the semantic content of the underlay's success rule: when
\(\alpha\) removes \(n\), every vertex of its snapshot that is above \(n\) is
already garbage.  It is the only fact about removals that the linearization
argument uses.

\section{Why the abstract pop must be placed retroactively}

\begin{example}[The lagging push]\label{ex:f1}
Consider the following concrete trace, with clock values on the left.
\[
\begin{array}{cl}
0,1 & \text{thread 1 pushes } m \text{ and completes};\\
2 & \text{thread 2 starts pushing } n\ (n\to_\conc m);\\
3 & \text{thread 3 snapshots } \{m,n\};\\
4 & \text{thread 4 starts pushing } x\ (x\to_\conc m,\ x\not\to_\conc n);\\
5 & \text{thread 2 completes } n;\\
6 & \text{thread 5 snapshots } \{m,n,x\};\\
7 & \text{thread 5 removes } n;\\
8 & \text{thread 3 removes } m\ (\text{top of }\{m,n\}\setminus\{n\}).
\end{array}
\]
The final removal is legal for TryStackAux: \(x\) is not in thread~3's
snapshot.  For the atomic TryStack, however, \(m\) can only be popped when
no live vertex is above it.  Both \(n\) and \(x\) are above \(m\)
concretely.  The pop of \(m\) can be linearized right after thread~3's
snapshot at clock~3 \emph{provided the abstract push-inv of \(n\) is placed
after it}: abstractly, \(n\)'s push then starts at ``3.1'' even though it
started at~2 concretely.  This is sound because \(n\)'s push is still pending
at clock~3.  Thread~5's pop of \(n\) at clock~6 is then legal because
\(x\)'s abstract push-inv (clock~4) precedes \(n\)'s response (clock~5), so
\(x\) is not above \(n\) abstractly.
\end{example}

The example shows that the abstract push-inv of a vertex may lag behind its
concrete push-inv, and that the correct lag depends on events that happen
later.  The proof therefore does not commit to a particular placement; it
keeps every placement that is consistent with the events recorded so far.

\section{Abstract histories}

\subsection{Positions}

Abstract events are placed into the gaps between concrete events.  A position
is a pair
\[
  p=(g,k)\in\Pos\triangleq\mathbb N\times\mathbb N,
\]
meaning ``in the gap right after the concrete event with clock value \(g\),
at rank \(k\)''.  Positions are ordered lexicographically, written
\(p\lt q\).  The rank budget of a gap is bounded,
\[
  \mathsf{ok}(p)\ \triangleq\ \rank(p)<2\gap(p)+2,
\]
and the rank \(2g+2\) is reserved for the abstract push response of the
concrete event at \(g\):
\[
  \ret(r)\triangleq(r,2r+2).
\]
Thus the abstract response of a completed push is fixed at its concrete
response (it is the last event of its gap), while abstract push-invs and pops
float within the budgeted ranks.  The bound is what makes the replay argument
of Section~\ref{sec:replay} a finite iteration over ranks.

\subsection{Histories and consistency}

\begin{definition}[History]
A history is a pair \(h=(h_\inv,h_\pop)\) with
\(h_\inv:\Node\rightharpoonup\Pos\) (the abstract push-inv of each vertex)
and \(h_\pop:\Node\rightharpoonup\Pos\times\Tid\) (the abstract pop of a
vertex and the popping thread).
\end{definition}

Write \(x\to_h y\) for the \emph{abstract edge}: \(h_\inv(x)=p\), \(r_y\)
is defined and \(\ret(r_y)\lt p\), i.e.\ the abstract push of \(x\) starts
after the response of \(y\).

\begin{definition}[Consistency, \code{consistent}]\label{def:consistent}
A history \(h\) is consistent with a ghost-timed state \(\conc\) when:
\begin{enumerate}[nosep,label=(C\arabic*)]
\item\label{c:inv} if \(h_\inv(n)=p\) then \(\mathsf{ok}(p)\), \(\gap(p)<\now\),
  \(i_n\le\gap(p)\), and \(\gap(p)\le r_n\) whenever \(r_n\) is defined:
  the abstract push-inv lies between the concrete push-inv and the concrete
  response;
\item\label{c:ret} if \(r_n\) is defined then \(h_\inv(n)\) is defined;
\item\label{c:pop} if \(h_\pop(n)=(p,\alpha)\) then \(\mathsf{ok}(p)\),
  \(\gap(p)<\now\), \(h_\inv(n)=q\lt p\), and either
  \(\rem(n)=(\alpha,st,rt)\) with \(st\le\gap(p)\le rt\)
  (a \emph{committed} pop, placed within the remover's interval), or
  \(\rem(n)\) is undefined, \(\snap(\alpha)=st\le\gap(p)\) and \(i_n<st\)
  (a \emph{speculative} pop of a vertex in \(\alpha\)'s pending snapshot);
\item\label{c:remp} every removed vertex is popped by its remover;
\item\label{c:spec} a thread's pending snapshot speculates on at most one
  vertex;
\item\label{c:top} (\emph{top}) if \(h_\pop(n)=(p,\alpha)\) and \(x\) has
  \(h_\inv(x)=q\lt p\) and \(x\to_h n\), then \(h_\pop(x)=(p',\_)\) with
  \(p'\lt p\): every abstractly newer vertex whose push started before the
  pop was already popped;
\item\label{c:uniq} all positions used by \(h\) are pairwise distinct.
\end{enumerate}
\end{definition}

Condition~\ref{c:top} is exactly the TryStack pop rule read along the
abstract order: a vertex is popped only while it is a top of the abstract
live graph.  Conditions~\ref{c:inv} and~\ref{c:pop} bound the retroactive
freedom by the concrete intervals of the operations involved, which is what
linearizability permits.

\subsection{Phases}

The proof needs to know, for each active thread, which operation it invoked
and how far it has progressed.  A \emph{phase map} \(\varphi:\Tid\rightharpoonup\)
\[
  \{\phase{PushInvoked}(v),\ \phase{PushPending}(v,\ell),\ \phase{PushDone}(v),\
    \phase{PopInvoked},\ \phase{PopSnapshot},\ \phase{PopEmptyPending},\
    \phase{PopDone}(r)\}
\]
records this.  It is \emph{consistent} with a control state \(c\) (payload
\(\conc\)) when \(P(t)=\ell\) iff \(\varphi(t)=\phase{PushPending}(\_,\ell)\);
\(S(t)\) is defined iff \(\varphi(t)=\phase{PopSnapshot}\);
\(\varphi(t)=\phase{PopEmptyPending}\) iff \(c=\Atomic(\conc,t,\TryPop)\);
a pending push's node carries its value; a \(\phase{PushDone}(v)\) thread has a
latest completed node of value \(v\); a \(\phase{PopDone}(\Succ(v,\beta,\ell))\)
thread is the recorded remover of \((\beta,\ell)\), which has value \(v\);
and \(\phase{PopEmptyPending}\) implies that every vertex is garbage.

\section{Representation of an abstract configuration}

\subsection{Cuts and tokens}

A history describes all abstract events; the abstract state \emph{at a
position} is obtained by executing the events before it.  For a position
\(C\) (the \emph{cut}) define
\begin{align*}
  \Placed_h(C,n)&\triangleq h_\inv(n)=q\lt C,\\
  \Popped_h(C,n)&\triangleq h_\pop(n)=(q,\_)\lt C,\qquad
  \Popped_h(C,\alpha,n)\ \text{with the popper fixed},\\
  \RetBefore_\conc(C,n)&\triangleq r_n\ \text{defined and}\ \ret(r_n)\lt C.
\end{align*}
The framework represents the progress of thread \(t\)'s current operation
\(m\) by a token \(\tokens(t)\in\{\Inv(m),\Lini(m),\Linr(m,r)\}\): invoked
but not yet linearized, interval-invocation linearized, response linearized
with result \(r\).  The token that thread \(t\) must carry at cut \(C\) is
determined by its phase and the history:

\begin{center}
\renewcommand{\arraystretch}{1.25}
\begin{tabular}{@{}ll@{}}
\toprule
\(\varphi(t)\) & \(\Tok_{\conc,\varphi,h}(C,t)\)\\
\midrule
undefined & no token\\
\(\phase{PushInvoked}(v)\) & \(\Inv(\Push v)\)\\
\(\phase{PushPending}(v,\ell)\) &
  \(\Lini(\Push v)\) if \(\Placed(C,(t,\ell))\), else \(\Inv(\Push v)\)\\
\(\phase{PushDone}(v)\), latest node \((t,\ell)\) &
  \(\Linr(\Push v,())\) if \(\RetBefore(C,(t,\ell))\); else
  \(\Lini\) if placed, else \(\Inv\)\\
\(\phase{PopInvoked}\), \(\phase{PopEmptyPending}\) & \(\Inv(\TryPop)\)\\
\(\phase{PopSnapshot}\) &
  \(\Linr(\TryPop,\Succ(V(n),n))\) if \(t\) speculatively popped a live
  \(n\) before \(C\); else \(\Inv(\TryPop)\)\\
\(\phase{PopDone}(\Succ(v,\beta,\ell))\) &
  \(\Linr(\TryPop,\Succ(v,\beta,\ell))\) if \(\Popped(C,t,(\beta,\ell))\),
  else \(\Inv(\TryPop)\)\\
\(\phase{PopDone}(\Fail)\), \(\phase{PopDone}(\Succ\Empty)\) &
  \(\Linr(\TryPop,r)\) (or \(\Inv(\TryPop)\) for the one thread whose
  response is being replayed)\\
\bottomrule
\end{tabular}
\end{center}

The last row uses a parameter \(pd\) (``pending done'') that names the
thread whose failed or empty response is currently being replayed; outside
the replay it is \(\bot\).

\begin{definition}[Representation at a cut, \code{represents_at}]
\(\Rep_{\conc,\varphi,h}(C,pd,\abst,\tokens)\) holds when
\(\abst=\Ready(V',E',P',g')\) and
\begin{align*}
  V'(n)&=V(n) &&\text{if }\Placed(C,n),\qquad V'(n)=\bot\ \text{otherwise};\\
  E'(x,y)&\iff \Placed(C,x)\land x\to_h y;\\
  P'(t)=\ell&\iff \Placed(C,(t,\ell))\land\lnot\RetBefore(C,(t,\ell));\\
  g'(n)&\iff\Popped(C,n);\\
  \tokens(t)&=\Tok_{\conc,\varphi,h}(C,t)\quad\text{for all }t.
\end{align*}
\end{definition}

The abstract graph is thus the concrete vertex map restricted to the vertices
whose abstract push has started, with abstract edges, pending pushes that
have started abstractly but not returned, and garbage equal to the popped
vertices.

\subsection{Realizable configurations}

The \emph{final cut} of a state is \(\cut(\conc)\triangleq(\now,0)\): the
position right after the last concrete event.  The abstract configurations
realizable from control state \(c\) with phases \(\varphi\) are
\[
  \Real(c,\varphi)\ \triangleq\
  \{(\abst,\tokens)\mid \exists h.\ h\text{ consistent with }\conc\land
     \Rep_{\conc,\varphi,h}(\cut(\conc),\bot,\abst,\tokens)\},
  \qquad \conc=\mathrm{payload}(c).
\]

\section{The invariant}

A proof state is \(w=(c,\poss)\) where \(c\) is the concrete control state
and \(\poss\) is a nonempty set of possibilities \((\abst,\tokens)\) sharing
an active-thread domain.

\begin{invariant}[\(J\) and \(I\)]
\(J(c,\varphi,\poss)\) is the conjunction of
\begin{enumerate}[nosep]
\item ghost well-formedness of \(c\) (Invariant~\ref{inv:wf});
\item phase consistency of \(\varphi\) with \(c\);
\item soundness: \(\poss\subseteq\Real(c,\varphi)\);
\item completeness: for every history \(h\) consistent with \(\conc\) there
  is \((\abst,\tokens)\in\poss\) with
  \(\Rep_{\conc,\varphi,h}(\cut(\conc),\bot,\abst,\tokens)\).
\end{enumerate}
The global invariant is \(I(c,\poss)\triangleq\exists\varphi.\ J(c,\varphi,\poss)\).
\end{invariant}

Soundness says every retained possibility is explained by a consistent
history; completeness says every consistent history is represented.  The
pair pins \(\poss\) down to \(\Real(c,\varphi)\) up to the choice of
representatives, and completeness is what later lets the proof select, after
a concrete response, the branch that realizes the observed result.

\subsection{Phases are determined by tokens}

\begin{lemma}[Token shape, \code{J_token_shape}]\label{lem:shape}
If \(J(c,\varphi,\poss)\) and \((\abst,\tokens)\in\poss\), then for every
thread \(t\) the token \(\tokens(t)\) has the shape prescribed by
\(\varphi(t)\): no token iff no phase; \(\Inv(\Push v)\) with \(P(t)\)
undefined for \(\phase{PushInvoked}(v)\); \(\Inv\) or \(\Lini\) of
\(\Push v\) with \(P(t)=\ell\) for \(\phase{PushPending}(v,\ell)\);
\(\Linr(\Push v,())\) for \(\phase{PushDone}(v)\); \(\Inv(\TryPop)\) with no
snapshot for \(\phase{PopInvoked}\); \(\Inv(\TryPop)\) or a
\(\Linr(\TryPop,\Succ(\ldots))\) with a snapshot for \(\phase{PopSnapshot}\);
\(\Inv(\TryPop)\) in the atomic-pending control state for
\(\phase{PopEmptyPending}\); \(\Linr(\TryPop,r)\) for \(\phase{PopDone}(r)\).
\end{lemma}

The proof unfolds \(\Tok\) at the final cut.  Two cases need
Invariant~\ref{inv:wf}: for \(\phase{PushDone}(v)\) the latest node has a
response \(r<\now\), so \(\ret(r)\lt\cut(\conc)\) and the token is
\(\Linr\); for \(\phase{PopDone}(\Succ(v,\beta,\ell))\) the removal is
recorded, so by~\ref{c:remp} and~\ref{c:pop} its pop lies below the final
cut and the token is \(\Linr\).

\begin{lemma}[Phase determinacy, \code{phase_from_token}, \code{phase_unique}]
\label{lem:determinacy}
If \(J(c,\varphi,\poss)\), \(J(c,\varphi',\poss')\), \((\abst,\tokens)\in\poss\),
\((\abst',\tokens')\in\poss'\) and \(\tokens(t)=\tokens'(t)\), then
\(\varphi(t)=\varphi'(t)\).  In particular any two phase maps witnessing
\(I\) for the same proof state agree everywhere.
\end{lemma}

The shapes of Lemma~\ref{lem:shape} are pairwise incompatible once the
concrete control state is fixed: pending-push entries, snapshots and the
atomic control state separate the push phases from each other and the pop
phases from each other, and \(\Linr\) results are injective.  Consequently
\(\varphi\) is a function of the proof state, and assertions may quantify
over ``the'' phase map.

\section{Rely and guarantee}

\begin{definition}
\begin{align*}
  G_\alpha(w,w')&\triangleq I(w)\land I(w')\land
    \forall\varphi,\varphi'.\ J\text{-witnesses for }w,w'\Rightarrow
    \forall t\ne\alpha.\ \varphi'(t)=\varphi(t),\\
  R_\beta(w,w')&\triangleq I(w)\land
    \forall\varphi,\varphi'.\ J\text{-witnesses for }w,w'\Rightarrow
    \varphi'(\beta)=\varphi(\beta).
\end{align*}
\end{definition}

A step of \(\alpha\) may change anything about the concrete state and the
possibility set, as long as the invariant holds and no other thread's phase
changes.  Correspondingly, an observer relies only on its own phase being
preserved.  This is enough because every method assertion below is of the
form ``\(I\) holds and my phase satisfies a predicate''.

\begin{lemma}[Validity and compatibility]
\begin{enumerate}[nosep]
\item (\code{valid_rg}) Under \(R_\beta\), thread \(\beta\) has a token in
  all possibilities of \(w\) iff it has one in all possibilities of \(w'\).
  By Lemma~\ref{lem:shape}, having a token is equivalent to having a phase,
  and \(R_\beta\) preserves the phase.
\item (\code{parallel_compatible}) For \(\alpha\ne\beta\), each of
  \(G_\alpha\), the framework's invocation step of \(\alpha\) (which adds
  \(\Inv(m)\) to \(\alpha\)'s token in every possibility), its return step
  (which removes \(\Linr(m,r)\)), and the identity, implies \(R_\beta\).
  For the invocation and return steps, the token of \(\beta\) is unchanged,
  so Lemma~\ref{lem:determinacy} gives \(\varphi'(\beta)=\varphi(\beta)\).
\end{enumerate}
\end{lemma}

The framework's invocation and return steps also preserve \(J\)
(\code{J_ac_inv}, \code{J_ac_res}): the phase map is updated at \(\alpha\)
to \(\phase{PushInvoked}(v)\) or \(\phase{PopInvoked}\), respectively removed,
and the representation of every history is adjusted by the same token
update, which matches the \(\Tok\) table.

\section{Simulating a concrete event}\label{sec:replay}

\subsection{The generic step}

Every concrete event is simulated by one lemma scheme.  Let \(c\) step to
\(c'\), with payloads \(\conc,\conc'\), and let \(\varphi'\) be the phase map
after the event (the acting thread's phase changes, all others are kept).

\begin{lemma}[\code{real_step_generic}]\label{lem:generic}
Assume \(J(c,\varphi,\poss)\), ghost well-formedness of \(c'\), phase
consistency of \(\varphi'\) with \(c'\), and the \emph{event property}:
for every history \(h'\) consistent with \(\conc'\),
\begin{enumerate}[nosep]
\item[(a)] the restriction \(h'{\restriction}\cut(\conc)\) (the events of
  \(h'\) below the old final cut) is consistent with \(\conc\);
\item[(b)] every possibility representing
  \(h'{\restriction}\cut(\conc)\) at \(\cut(\conc)\) for \((\conc,\varphi)\)
  can take finitely many abstract steps to a possibility representing
  \(h'\) at \(\cut(\conc')\) for \((\conc',\varphi')\).
\end{enumerate}
Then there is \(\poss'\), contained in the step-closure of \(\poss\), with
\(J(c',\varphi',\poss')\).
\end{lemma}

\begin{proof}
Take
\[
  \poss'\triangleq\{(\abst',\tokens')\mid
    (\abst',\tokens')\text{ reachable from }\poss\land
    (\abst',\tokens')\in\Real(c',\varphi')\}.
\]
Soundness of \(\poss'\) is immediate.  For completeness, let \(h'\) be
consistent with \(\conc'\).  By (a) and the completeness of \(\poss\), some
\((\abst,\tokens)\in\poss\) represents \(h'{\restriction}\cut(\conc)\); by
(b) it steps to a representative of \(h'\), which lies in \(\poss'\).
Finally \(\poss'\) is nonempty because the trace theorem
(Section~\ref{sec:trace}) supplies a consistent history of \(\conc'\), and
completeness applied to it exhibits a member.
\end{proof}

The nonemptiness step is where the trace theorem is indispensable: the
framework's possibility sets are intrinsically nonempty, so a concrete
event whose every explanation had been discarded would be a genuine proof
failure.

\subsection{What the two parts of the event property mean}

A history is a choice of linearization points: condition~\ref{c:inv} places
the abstract push-inv of \(n\) inside the real-time interval of its push,
\ref{c:pop} places the pop of \(n\) inside the interval of the try-pop that
removed it (or, speculatively, after the snapshot of a try-pop still in
progress), and the remaining conditions say that executing the events in
position order is a legal run of the atomic TryStack.  A consistent history
of the current state is therefore precisely a \emph{linearization of the
trace so far}, and a possibility \((\abst,\tokens)\in\poss\) is the abstract
state reached by running one such linearization up to the present moment.
The invariant keeps all of them because the trace so far does not determine
which linearization the future will need.

Part~(a), \emph{restriction}, says that every linearization of the longer
trace, when cut back to the events already placed before the present
moment, is a linearization of the shorter trace.  Semantically this is the
statement that the new concrete event never invalidates an already
observable abstract past: the abstract events below the old cut form a run
that was already legal, and the new event only adds abstract events in the
gap just opened.  Part~(a) is what allows the proof to be a
\emph{forward} simulation over a set: it guarantees that the explanations
of the extended trace are extensions of explanations of the prefix, so the
family of possibilities carried so far contains a representative of every
future that remains possible.  Without it, an event could require an
abstract past that no retained possibility has, and completeness would be
lost.  Note the quantification: it ranges over histories of the \emph{new}
state.  A history of the old state need not extend, and the proof never
claims it does.  Linearizations of the prefix that no longer extend are
simply dropped, which is the set-valued analogue of a backward simulation
or a prophecy variable choosing the right future.

Part~(b), \emph{replay}, says that the abstract machine can actually
execute the events that the new linearization places in the fresh gap,
starting from a possibility that has executed the events before it.  This
is the half that connects the invariant to the layer simulation: the
framework demands that possibilities evolve only by abstract steps, because
those steps are exactly what shows that the specification \(\mathcal F\)
accepts the linearized trace.  Consistency conditions are stated on
positions and clock values; part~(b) turns them into transitions of the
TryStack LTS: the top condition~\ref{c:top} becomes the enabledness of the
atomic pop, the interval conditions become the presence or absence of the
pending-push entry, and the token table becomes the update of the
linearization tokens that the framework reads at method return.

Together the two parts are the induction step of a linearizability proof
phrased over sets of partial linearizations.  Linearizability asks for one
total order of the operations of the complete trace that respects real time
and is accepted by the specification.  Part~(a) ensures that such an order,
whenever it exists for the complete trace, is visible at every prefix as a
retained possibility; part~(b) ensures that the specification accepts each
prefix of that order, step by step; and the trace theorem
(Section~\ref{sec:trace}) ensures that at least one such order exists for
every reachable state, so that the retained set is never empty.  At a method
return the framework then reads the acting thread's token, which by
Lemma~\ref{lem:shape} is \(\Linr(m,r)\) in every retained possibility, and
concludes that every possible linearization returns the observed result.

\subsection{Restriction}

Part (a) of the event property is proved once for all events
(\code{restrict_consistent}).  Let \(\conc\) and \(\conc'\) be related by
\emph{ghost extension}: clock values recorded in \(\conc'\) below
\(\now(\conc)\) are exactly those of \(\conc\), a removal recorded in
\(\conc'\) with \(rt\ge\now(\conc)\) corresponds in \(\conc\) to a pending
snapshot at the same \(st\), and so on.  Every step of TryStackAux is a ghost
extension of its pre-state.  Then for \(h'\) consistent with \(\conc'\), the
history \(h'{\restriction}\cut(\conc)\) is consistent with \(\conc\): the
push-inv and pop conditions transfer because the relevant clock values are
below \(\now(\conc)\); a committed pop of \(h'\) whose removal is not yet
recorded in \(\conc\) becomes a \emph{speculative} pop in \(\conc\), which
is legitimate because the popper's snapshot is pending in \(\conc\) with the
same snapshot time; and the top condition transfers because abstract edges
below the cut are the same.  A side condition needed for~\ref{c:spec} is
that two committed pops by the same thread cannot both fall below the cut
while being unrecorded in \(\conc\); it holds because a thread has one
snapshot at a time.

\subsection{Replay}

Part (b) is a replay of the abstract events of \(h'\) that lie between the
two cuts.  One position at a time, the representation is moved from cut
\(C\) to \(\mathsf{next}(C)\) (\(C\) with the rank incremented):
\begin{itemize}[nosep]
\item If \(h'_\inv(n)=C\) (\code{replay_inv}): the abstract state takes the
  TryStack push-invocation step of \(n=(t,\ell)\).  The vertex \(n\) is
  fresh in the represented state (not placed before \(C\), not popped),
  the edges it gains are exactly the abstract edges \(n\to_{h'}y\) with \(y\)
  placed before \(C\) (a vertex placed and non-pending at \(C\) is one
  whose response precedes \(C\), and \(\ret(r_y)\lt C=h'_\inv(n)\) is the
  definition of \(n\to_{h'}y\)), and thread \(t\)'s token becomes
  \(\Lini(\Push v)\), as the \(\Tok\) table requires once \(n\) is placed.
\item If \(h'_\pop(n)=(C,\alpha)\) (\code{replay_pop}): the abstract state
  takes the atomic try-pop invocation and success response of \(\alpha\)
  for \(n\).  Condition~\ref{c:top} shows that \(n\) is a top of the
  represented live graph: any live \(x\) with an abstract edge \(x\to n\) is
  placed before \(C\) and therefore popped before \(C\), i.e.\ garbage.  The
  token of \(\alpha\) becomes \(\Linr(\TryPop,\Succ(V(n),n))\).
\item If \(C=\ret(r)\) for the response of \((t,\ell)\) (\code{replay_ret}):
  the abstract push-response step of \(t\) is taken; the pending entry
  \(P'(t)\) is removed and the token becomes \(\Linr(\Push v,())\).  The
  entry is present because~\ref{c:ret} and~\ref{c:inv} place \(h'_\inv(t,\ell)\)
  before \(\ret(r)\).
\item Otherwise nothing happens at \(C\) and the representation is
  unchanged (\code{represents_at_next_none}).
\end{itemize}
In each case the tokens of all other threads are unchanged because the
predicates \(\Placed\), \(\Popped\), \(\RetBefore\) change only at the
thread whose event sits at \(C\).  Iterating from rank \(0\) to \(2g+3\)
replays a whole gap (\code{replay_gap}); by the rank budget the cut
\((g,2g+3)\) sees the same events as \((g{+}1,0)\) (\code{cut_equiv_gap_end}),
so one gap of replay moves the final cut forward by one clock tick
(\code{replay_last_gap}).  Failed and empty pop responses are replayed
separately (\code{replay_done}): the two abstract events are adjacent, and
for the empty result the side condition ``every placed vertex is popped''
is what the TryStack empty rule requires.

The phases needed by the replay (\code{gap_phases}) hold in the last gap:
a push-inv placed in the gap of the concrete event just recorded belongs
either to the pending push of the acting thread or to a completed push
(whose thread is \(\phase{PushDone}\)); a pop placed there is either the
committed pop of the acting thread or a speculative pop of a thread in
\(\phase{PopSnapshot}\).

\subsection{The event lemmas}

For each of the seven TryStackAux transitions there is a lemma
(\code{real_step_push_inv}, \code{real_step_push_res},
\code{real_step_snapshot}, \code{real_step_remove}, \code{real_step_fail},
\code{real_step_empty_inv}, \code{real_step_empty_res}) establishing the
event property with the appropriate phase update:

\begin{center}
\renewcommand{\arraystretch}{1.2}
\begin{tabular}{@{}lll@{}}
\toprule
concrete event of \(\alpha\) & new phase of \(\alpha\) & abstract events replayed\\
\midrule
\(\Push(v)\) invocation at \(\ell\) & \(\phase{PushPending}(v,\ell)\) &
  the events of the new last gap (possibly \(h'_\inv(\alpha,\ell)\))\\
\(\Push\) response & \(\phase{PushDone}(v)\) &
  the last gap, ending with \(\ret\) of \((\alpha,\ell)\)\\
snapshot invocation & \(\phase{PopSnapshot}\) & the last gap\\
success response for \(n\) & \(\phase{PopDone}(\Succ(V(n),n))\) &
  the last gap; the committed pop of \(n\) lies below the new cut\\
failure response & \(\phase{PopDone}(\Fail)\) &
  the last gap, then the abstract \(\Fail\) pop of \(\alpha\)\\
empty invocation & \(\phase{PopEmptyPending}\) & none (control state only)\\
empty response & \(\phase{PopDone}(\Succ\Empty)\) &
  the abstract \(\Succ(\Empty)\) pop of \(\alpha\)\\
\bottomrule
\end{tabular}
\end{center}

The interesting transfer facts are the following.
\begin{itemize}[nosep]
\item A push invocation adds a vertex that no history of the pre-state can
  mention; the token of \(\alpha\) is \(\Inv(\Push v)\) before and after the
  old cut, matching \(\phase{PushPending}\) with the node not yet placed.
\item A push response makes the responded node \((\alpha,\ell)\) the latest
  node of \(\alpha\).  Its abstract push-inv, if it had not been placed
  already, must be placed in the last gap; the replay places it and then
  takes the response.
\item A success response for \(n\) must, in every retained branch, have
  popped \(n\) before the new cut: \ref{c:remp} and~\ref{c:pop} give
  \(h'_\pop(n)=(p,\alpha)\) with \(\gap(p)\le rt=\now(\conc)\).  Branches of
  the pre-state in which \(\alpha\) had speculatively popped a different
  vertex, or none, are those whose history does not restrict from any
  consistent history of the post-state; they are dropped by the definition
  of \(\poss'\).  The concrete garbage update is mirrored abstractly because
  the represented garbage is the popped set.
\item A failure response is possible in every branch in which \(\alpha\) has
  not speculatively popped: the token \(\Inv(\TryPop)\) is turned into
  \(\Linr(\TryPop,\Fail)\) by the adjacent abstract invocation and failure
  response.  Branches with a speculative pop are dropped.
\item The empty branch is atomic on both sides.  At the concrete invocation
  the phase becomes \(\phase{PopEmptyPending}\) with the token unchanged; at
  the response the abstract events are replayed with the ``pending done''
  parameter set to \(\alpha\).  The TryStack empty rule needs every abstract
  vertex to be garbage; the concrete side condition says every concrete
  vertex is garbage, i.e.\ removed, hence popped below the cut by
  \ref{c:remp}.
\end{itemize}

Phase consistency of \(\varphi'\) with \(c'\) is checked separately for
each event (\code{pc_push_inv}, \ldots, \code{pc_empty_res}); ghost
well-formedness of \(c'\) is the preservation theorem of
Invariant~\ref{inv:wf}.

\section{The method triples}

\subsection{Assertions}

All assertions are of the form
\[
  \PhaseIs_\alpha(Q)\ \triangleq\ I\land\forall\varphi.\ J\text{-witness}\Rightarrow Q(\varphi(\alpha)),
\]
which by Lemma~\ref{lem:determinacy} is a statement about the unique phase
of \(\alpha\).  It entails \(I\) and is stable under \(R_\alpha\), since the
rely preserves \(\alpha\)'s phase.  The postcondition of an operation is
\[
  \Done_\alpha(m,r)\triangleq
  \PhaseIs_\alpha(\phase{PushDone}(v))\ \text{for }m=\Push v,\qquad
  \PhaseIs_\alpha(\phase{PopDone}(r))\ \text{for }m=\TryPop.
\]
By Lemma~\ref{lem:shape}, \(\Done_\alpha(m,r)\) implies that every
possibility carries the token \(\Linr(m,r)\) for \(\alpha\), which is the
framework's requirement at a method return.

The framework's invocation step exposes the precondition: after the
invocation of \(m\) by \(\alpha\) the assertion
\(\PhaseIs_\alpha(\phase{PushInvoked}(v))\), respectively
\(\PhaseIs_\alpha(\phase{PopInvoked})\), holds (\code{ginv_exposes_phase});
after the return step from \(\Done_\alpha(m,r)\) the invariant \(I\) holds
again with \(\alpha\)'s phase removed (\code{gret_closes_done}).

\subsection{Threads outside the domain}

If \(\alpha\notin\Upsilon\), both layers raise an error at the invocation.
The precondition is split (\code{phase_or_error}) into
\(\Inside_\alpha(P)\triangleq P\land\alpha\in\Upsilon\) and the abstract
error assertion: from a possibility with token \(\Inv(m)\) and
\(\abst=\Ready(\_)\) (the represented state is always \(\Ready\)), the
TryStack error transition is reachable.  Inside the domain the underlay
invocation raises no error.

\subsection{Push}

\[
  \{\PhaseIs_\alpha(\phase{PushInvoked}(v))\}\quad
  \Push_\alpha(v)\quad
  \{r.\ \Done_\alpha(\Push v,r)\}.
\]
The underlay call is a single visible operation.  Its invocation update
(\code{push_inv_update}) moves the phase to \(\phase{PushPending}(v,\ell)\)
for the location \(\ell\) chosen by the underlay and is justified by
Lemma~\ref{lem:generic} with the push-invocation event lemma.  The location
is recovered from the pre-state's pending map when the response update
(\code{push_res_update}) moves the phase to \(\phase{PushDone}(v)\).  Both
updates satisfy \(G_\alpha\) because only \(\alpha\)'s phase changes
(\code{G_intro}).  The trailing \(\mathsf{return}\) needs only stability of
the postcondition.

\subsection{Try-pop}

\[
  \{\PhaseIs_\alpha(\phase{PopInvoked})\}\quad
  \TryPop_\alpha\quad
  \{r.\ \Done_\alpha(\TryPop,r)\}.
\]
The invocation update (\code{trypop_inv_update}) inspects the concrete step:
the snapshot branch yields \(\phase{PopSnapshot}\), the empty branch
\(\phase{PopEmptyPending}\); the intermediate assertion is the disjunction.
The response update (\code{trypop_res_update}) again inspects the step.  If
the pre-state is atomic-pending the response is \(\Succ(\Empty)\) and the
phase was \(\phase{PopEmptyPending}\).  Otherwise the phase was
\(\phase{PopSnapshot}\) and the response is either a success for a vertex
\(n\) that is top in \(S(\alpha)\setminus g\) (phase
\(\phase{PopDone}(\Succ(V(n),n))\)) or a failure (phase
\(\phase{PopDone}(\Fail)\)).  In every case Lemma~\ref{lem:generic} with the
matching event lemma yields the new possibility set.

\subsection{Initial state and assembly}

In the initial state the clock is \(0\), no vertex exists, and the only
consistent history is empty; the singleton configuration with the initial
TryStack state and no tokens both is realizable and represents every
consistent history (\code{initial_I}).  The soundness theorem of the set
logic then combines validity of the rely--guarantee pair, parallel
compatibility, the two method triples with their exposure and closing
lemmas, and the initial invariant into \(\mathsf{MTryStack}\).

\section{The trace theorem}\label{sec:trace}

\begin{theorem}[\code{TryStackTrace.trace_theorem}]
Every ghost-well-formed TryStackAux state admits a consistent history.
\end{theorem}

This is the only fact about the specifications that the layer proof does
not obtain from the event lemmas.  It states that the retroactive placement
that the invariant relies on always exists, and it is a genuinely global
combinatorial statement: Example~\ref{ex:f1} shows that the right placement
of an early event depends on late events.

\subsection{Candidate placements that do not work}

Several natural placement rules fail, which is why the proof below is
inductive and chooses positions by a least-fixpoint criterion.  Write
\(\rem(n)=(\alpha_n,st_n,rt_n)\) for a removed vertex.
\begin{itemize}[nosep]
\item \emph{Pop at the snapshot, pushes lock-step.}  Let \(x\to_\conc n\)
  with both completed, \(\alpha\) snapshot \(\{n,x\}\), then \(\beta\)
  snapshot the same set, \(\beta\) remove \(x\), \(\alpha\) remove \(n\).
  Then \(\pop(x)\) is placed after \(\pop(n)\) while \(x\) is above \(n\)
  and live, violating~\ref{c:top}.  The pop of \(n\) must wait for the pop
  of \(x\).
\item \emph{Pop at the removal.}  Let \(\alpha\) snapshot \(\{n\}\), then a
  push of \(x\) start and complete (\(x\to_\conc n\)), then \(\alpha\)
  remove \(n\).  The pop of \(n\) must precede the abstract push-inv of
  \(x\), which cannot be delayed past \(x\)'s completion; the pop must be
  placed before \(i_x\).
\item \emph{Wait for all blockers.}  Combining the two shapes, waiting for
  every above-\(n\) vertex of the snapshot to be popped can push
  \(\pop(n)\) past the push of a completed, never-removed vertex above
  \(n\).  The correct rule postpones the push-inv of a blocker that is still
  pending instead of waiting for its pop.
\end{itemize}
These shapes were found by exhaustive search over short traces before the
mechanization; the construction below was validated the same way.

\subsection{Abstract histories with unbounded ranks}

The construction ignores the rank budget and restores it at the end.  An
\emph{abstract history} is a history whose positions satisfy the
consistency conditions of Definition~\ref{def:consistent} except
\(\mathsf{ok}\), with the convention that the push response of \(m\) follows
every abstract event of gap \(r_m\) (so \(x\to_h y\) simply means
\(r_y<\gap(h_\inv(x))\)), with all pops committed (no speculative pops), and
with two additional \emph{induction invariants}:
\begin{enumerate}[nosep,label=(I\arabic*)]
\item\label{i:tot} every vertex has an abstract push-inv;
\item\label{i:nd} (\emph{no needless delay}) if \(x\) is not popped,
  \(m\) is completed, \(i_x\le r_m<\gap(h_\inv(x))\), then \(m\) is popped
  and \(h_\pop(m)\lt h_\inv(x)\).
\end{enumerate}
\ref{i:nd} says that the abstract push-inv of an unpopped vertex is delayed
past the response of a vertex it does not concretely follow only when that
vertex was popped before the delayed position.  Without it, a history could
delay an unpopped \(x\) past \(\ret(r_m)\) for an \(m\) that is removed
later while \(x\) is still live, and then the removal of \(m\) would have no
consistent placement.

\subsection{Reverse induction on the clock}

The theorem is proved by induction on \(\now\).  For a well-formed state
with \(\now=T+1\), the event recorded at clock \(T\) is undone, producing a
well-formed state with clock \(T\) (Invariant~\ref{inv:wf} is re-established
by hand for each of the four undo operations; the push-response undo needs
clause~7 to show the pusher has no other pending node).  The abstract
history of the smaller state is then extended:
\begin{itemize}[nosep]
\item push invocation of \(x\) at \(T\): place \(h_\inv(x)=(T,0)\);
\item push response, snapshot, or no event at \(T\): keep the history;
\item removal of \(n\) at \(T\) with snapshot time \(st\): insert the pop of
  \(n\) as described next.
\end{itemize}

\subsection{Inserting a pop}

Fix the removal \(\rem(n)=(\alpha,st,T)\) and the abstract history \(h\) of
the state without it.  If \(n\) is pending, no vertex is abstractly above
\(n\) and the pop is placed at \((T,0)\).  Otherwise let \(r_n\) be its
response and call \(x\) \emph{newer} if \(r_n<\gap(h_\inv(x))\).

\begin{definition}[Feasible and safe positions]
For a position \(Q\), a \emph{blocker} is a newer vertex \(x\) with
\(h_\inv(x)\lt Q\) that is not popped before \(Q\).  The position \(Q\) is
\emph{feasible and safe}, \(\FS(Q)\), when
\begin{enumerate}[nosep]
\item \((st,0)\lte Q\), \(h_\inv(n)\lt Q\) and \(\gap(Q)\le T\);
\item every blocker \(x\) is \emph{pending at} \(Q\): \(r_x\ge\gap(Q)\) if
  defined;
\item every blocker \(x\) is \emph{safe}: for every popped \(m\) with
  \(\gap(h_\inv(x))\le r_m<\gap(Q)\) and \(h_\pop(m)\not\lt Q\), the vertex
  \(x\) is popped and \(h_\pop(x)\lt h_\pop(m)\).
\end{enumerate}
\end{definition}

Inserting the pop at \(P=(g,k)\) means: \(h_\pop(n)=(P,\alpha)\); every
existing event at gap \(g\) with rank at least \(k\) is shifted up by
\(T+2\); every blocker \(x\) has its push-inv moved to \((g,k+1+i_x)\),
right after the new pop and before the shifted events.  The moved vertices
are exactly those whose push-inv must be postponed behind the pop
(condition~2 makes that legal); condition~3 ensures that postponing does not
create an abstract edge \(x\to m\) for an \(m\) popped later while \(x\)
is live.

\begin{lemma}[Existence, \code{exists_M}]\label{lem:M}
Let \(C\) be the set of newer vertices that are completed and not popped.
If \(C=\emptyset\), then \(\FS((T,0))\).  Otherwise let \(M\) be the least
\(h_\inv(x_0)\) over \(x_0\in C\); then \(\FS(M)\).
\end{lemma}

\begin{proof}
For \(x_0\in C\): by \ref{i:nd} applied to \(x_0\) and \(n\) (which is not
popped in \(h\)), \(i_{x_0}>r_n\); by the blocking clause of
Invariant~\ref{inv:wf}, \(i_{x_0}\notin(r_n,st)\), for otherwise \(x_0\)
would be removed and hence popped by~\ref{c:remp}; and \(i_{x_0}\ne st\).
So \(i_{x_0}>st\), which gives \((st,0)\lt M\).  A blocker \(x\) at \(M\)
with \(r_x<\gap(M)\) is either in \(C\) with \(h_\inv(x)\lt M\), contradicting
minimality, or popped after \(M\); but \(x_0\) is newer than \(x\), placed
before that pop and unpopped, contradicting~\ref{c:top}.  Safety is
vacuous for the same reason: any \(m\) popped after \(M\) with
\(r_m<\gap(M)\) would see the unpopped \(x_0\) above it.
\end{proof}

Let \(P\) be the least position with \(\FS(P)\) (it exists classically);
by Lemma~\ref{lem:M}, \(P\lte M\).

\begin{lemma}[Moved vertices are already popped, \code{blocker_popped}]
\label{lem:U}
Every blocker of \(P\) is popped in \(h\) (necessarily after \(P\)).
\end{lemma}

\begin{proof}
Suppose \(x\) is an unpopped blocker.  As in Lemma~\ref{lem:M},
\(i_x>st\), so \(Q\triangleq h_\inv(x)\) satisfies condition~1 of \(\FS\).
Any blocker \(z\) of \(Q\) with \(r_z<\gap(Q)\) is either unpopped, hence in
\(C\) with \(h_\inv(z)\lt Q\lt P\lte M\), contradicting the minimality of
\(M\), or popped after \(Q\), contradicting~\ref{c:top} for the unpopped
newer \(x\).  Safety at \(Q\) is vacuous for the same reason.  Hence
\(\FS(Q)\) with \(Q\lt P\), contradicting the minimality of \(P\).
\end{proof}

\begin{lemma}[The extended history is an abstract history, \code{new_acons}]
\end{lemma}

\begin{proof}[Proof sketch]
The shift preserves the order of existing events and keeps them distinct
from the new positions.  The pop of \(n\) at \(P\) satisfies~\ref{c:pop}
since \(st\le g\le T\) and \(h_\inv(n)\lt P\) (\(n\) is not newer than
itself, so it is not moved).  For~\ref{c:top} at the new pop: a newer \(x\)
placed before \(P\) is not moved, hence not a blocker, hence popped before
\(P\).  For~\ref{c:top} at an old pop \(m\) shifted to \(p\): if \(x\) is
unmoved the claim is inherited; if \(x\) is moved and \(m\) is popped after
\(P\), then either \(x\) was already newer than \(m\) (inherited) or the
move crossed \(\ret(r_m)\), in which case safety gives
\(h_\pop(x)\lt h_\pop(m)\).  For~\ref{i:nd}: by Lemma~\ref{lem:U} moved
vertices are popped, so~\ref{i:nd} concerns only unmoved vertices, for
which the only new completed-and-popped pair is \((x,n)\), and
\ref{i:nd} of \(h\) excludes an unpopped \(x\) delayed past \(\ret(r_n)\)
since \(n\) was not popped in \(h\).  \ref{i:tot} is unchanged.
\end{proof}

\subsection{Restoring the rank budget}

An abstract history may use ranks beyond \(2g+1\).  Alongside the history
the induction maintains a \emph{key map} from clock values to abstract
events: the push-inv of \(x\) is keyed by \(i_x\) and the pop of \(n\) by
\(st_n\); both keys are at most the gap of the event by~\ref{c:inv}
and~\ref{c:pop}, and keys are distinct by clause~6 of
Invariant~\ref{inv:wf}.  Compression replaces the rank of an event at gap
\(g\) by the number of keys \(t\le g\) whose event lies at gap \(g\) with a
smaller rank.  This is an order isomorphism on the events of each gap, so
all consistency conditions transfer, and the new rank is at most \(g+1\),
within the budget \(2g+2\); in particular \(\ret(r)\) still follows every
compressed event of gap \(r\), so abstract edges are unchanged.  The key map
makes the count computable without any choice principle.

\section{Correspondence with the mechanized definitions}

\begin{center}
\renewcommand{\arraystretch}{1.15}
\begin{tabular}{@{}ll@{}}
\toprule
this note & Coq\\
\midrule
ghost fields \(\now,i_n,r_n,\snap,\rem\) &
  \code{tsa_now}, \code{tsa_node_inv}, \code{tsa_node_ret}, \code{tsa_snap_time}, \code{tsa_removals}\\
Invariant~\ref{inv:wf} & \code{TryStackAuxGhost.tsa_ghost_wf}, \code{step_preserves_wf}\\
positions, \(\mathsf{ok}\), \(\ret\) & \code{Pos}, \code{pos_lt}, \code{pos_ok}, \code{ret_pos}\\
history, \(x\to_h y\) & \code{History}, \code{abs_edge}\\
Definition~\ref{def:consistent} & \code{TryStackLinearization.consistent}\\
phases, phase consistency & \code{Phase}, \code{PhaseMap}, \code{phase_consistent}\\
\(\Tok\), \(\Rep\), \(\cut\), \(\Real\) &
  \code{token_at}, \code{represents_at}, \code{final_cut}, \code{Real}\\
\(J\), \(I\) & \code{TryStackProof.J}, \code{TryStackProof.I}\\
Lemmas~\ref{lem:shape}, \ref{lem:determinacy} &
  \code{J_token_shape}, \code{phase_from_token}, \code{phase_unique}\\
\(G_\alpha\), \(R_\beta\) & \code{G}, \code{R}, \code{valid_rg}, \code{parallel_compatible}\\
Lemma~\ref{lem:generic} & \code{real_step_generic}, \code{ac_real_step}\\
restriction & \code{restrict}, \code{ghost_extends}, \code{restrict_consistent}\\
replay & \code{replay_inv}, \code{replay_pop}, \code{replay_ret}, \code{replay_gap}, \code{replay_done}\\
event lemmas & \code{real_step_push_inv}, \ldots, \code{real_step_empty_res}\\
assertions & \code{PhaseIs}, \code{Done}, \code{Inside}, \code{phase_or_error}\\
triples & \code{push_method_triple}, \code{trypop_method_triple}\\
results & \code{MTryStack}, \code{MTryStackLinearizable}, \code{MListPoolTryStack}\\
abstract history, \ref{i:tot}, \ref{i:nd} & \code{TryStackTrace.acons} (\code{a_total}, \code{a_nd})\\
undo operations & \code{undo_inv}, \code{undo_ret}, \code{undo_snap}, \code{undo_rem}, \code{with_now}\\
\(\FS\), blockers, \(C\) & \code{FS}, \code{blocker}, \code{Cset}\\
Lemmas~\ref{lem:M}, \ref{lem:U} & \code{exists_M}, \code{blocker_popped}\\
insertion & \code{shift}, \code{moved_b}, \code{new_hist}, \code{new_acons}\\
compression & \code{KeyMap}, \code{cnt}, \code{compress}, \code{compress_consistent}\\
trace theorem & \code{TryStackTrace.trace_theorem}, used as \code{TryStackProof.real_nonempty}\\
\bottomrule
\end{tabular}
\end{center}

\end{document}
